\section{Results}
\label{sec:results}

\subsection{RQ1: Defect Patterns and Categories}
\label{sec:results-rq1}

We compare linter-detected violation patterns in AI-generated and human-authored code across six programming languages. AI-generated code exhibits significantly higher violation rates overall, with the magnitude of the gap depending strongly on language typing discipline and the specific violation category examined.

\paragraph{Overall violation density.}
A two-sample Mann--Whitney $U$ test on per-file violation counts confirms that AI-generated files contain significantly more violations than human-authored files ($U= \dots$, $p<0.001$). Cliff's $\delta \approx 0.58$ indicates a medium-to-large effect, meaning that in more than half of all pairwise comparisons a randomly selected AI-generated file accumulates more violations than a randomly selected human-authored file. This gap is not driven by a few extreme outlier files; the distributional shift is consistent across the full sample.

\paragraph{Severity profile.}
Higher violation counts would be less concerning if the excess violations were merely stylistic. To test this, we categorized all violations into three levels (Error, Warning, Info) and compared severity distributions with a chi-square test of independence. The two groups differ significantly ($\chi^2 = 24.73$, $df = 2$, $p < 0.001$, Cramér's $V = 0.18$). AI-generated code carries a substantially higher proportion of error-level violations (32.5\%) than human-authored code (18.7\%), while human code shows a larger share of warning-level issues (61.2\% vs.\ 48.3\% for AI). Error-level violations are those that can halt execution or indicate broken module dependencies; the near-twofold gap in error prevalence indicates that the elevated defect density in AI code is not cosmetic noise but includes a material share of functionally disruptive issues.

\paragraph{Language-level breakdown.}
Table~\ref{tab:violation_rates} reports per-file violation rates for each language. The differences are far from uniform. Python shows the largest gap: AI-generated files average 26.68 violations per file compared to 9.69 for human-authored files, a 2.75-fold increase. Python's dynamically typed and interpreted nature means that type mismatches, undefined names, and import failures that would be caught at compile time in other languages persist as linter warnings, leaving more room for generation-time errors. TypeScript, by contrast, shows only a 1.32-fold difference, consistent with the expectation that static type checking constrains the space of valid code a model can generate.

C\# presents an exception: AI-generated C\# files have \emph{fewer} violations per file (1.71) than human-authored files (3.93), a ratio of 0.43. C\#'s strict compilation semantics may force AI models toward syntactically valid patterns, and the human baseline drawn from active pull requests may include draft-quality files with higher stylistic variation. Java's human baseline records zero violations because those files were sourced from successfully merged pull requests subject to automated quality gates; the non-zero AI violation rate (4.99 violations per file) reflects the gap between curated human contributions and unconstrained model output rather than a property of Java in general.

Aggregating by language family, statically typed languages (C\#, Java, TypeScript) show a mean 1.66-fold gap, versus 2.34-fold for dynamically typed languages (Python, JavaScript). Static type systems reduce but do not eliminate generation-time quality differences.

\begin{table}[H]
\centering
\caption{Linter violation rates by language (V/F = violations per file).}
\label{tab:violation_rates}
\footnotesize
\resizebox{\columnwidth}{!}{%
\begin{tabular}{lrrrr}
\toprule
\textbf{Language} & \textbf{AI V/F} & \textbf{Human V/F} & \textbf{Ratio} & \textbf{Type} \\
\midrule
C\#         & 1.71  & 3.93 & 0.43 & Static  \\
C++         & 1.32  & 0.91 & 1.45 & Hybrid  \\
Java        & 4.99  & 0.00 & --   & Static  \\
JavaScript  & 2.86  & 2.92 & 0.98 & Dynamic \\
Python      & 26.68 & 9.69 & 2.75 & Dynamic \\
TypeScript  & 1.12  & 0.85 & 1.32 & Static  \\
\midrule
\textbf{Mean (Static)}  & \textbf{2.59}  & \textbf{1.56} & \textbf{1.66} & \\
\textbf{Mean (Dynamic)} & \textbf{14.77} & \textbf{6.31} & \textbf{2.34} & \\
\bottomrule
\end{tabular}%
}
\end{table}

\paragraph{Dominant violation categories.}
Table~\ref{tab:static_analysis_ai_vs_human} shows the two most frequently observed violations per language for each group. Several cross-language patterns emerge. Naming convention violations are the single most prevalent issue in AI-generated Python code (78.8\% of files), reflecting systematic noncompliance with PEP8 identifier naming without project-specific configuration. Import errors appear in 56.5\% of AI-generated Python files and 51.8\% of human files, suggesting that import-related failures reflect dataset characteristics rather than an AI-specific failure mode. Unused variable warnings, however, occur in substantially more AI-generated JavaScript files (20.6\%) than human files (8.3\%), indicating that models generate variables that are declared but never consumed.

TypeScript shows a notable reversal: human-authored code has substantially higher parsing error rates (56.9\% of files) than AI-generated code (14.2\%). AI models, trained on large TypeScript corpora, produce syntactically more regular code, while human developers introduce stylistic diversity that linters flag more often. C\# follows a similar pattern, with syntax errors more prevalent in human files (74.4\%) than AI files (47.3\%).

Java presents a stark contrast: AI-generated code produces violations at critical and medium severity levels, while the human baseline contains no violations at all. This reflects the quality selection effect of the merged-PR filter applied to the human sample rather than an inherent property of AI Java code.

\begin{table*}[t!]
\centering
\caption{Static Analysis Results Comparing AI-Generated and Human-Written Code Violations}
\label{tab:static_analysis_ai_vs_human}
\small
\begin{tabularx}{\textwidth}{@{}lllXXlXX@{}}
\toprule
\textbf{Language} & \textbf{Severity} & \textbf{AI Message} & \textbf{AI (\#)} & \textbf{AI (\%)} & \textbf{Human Message} & \textbf{Human (\#)} & \textbf{Human (\%)} \\
\midrule
\multirow{8}{*}{Python}
& \multirow{2}{*}{Error}
  & ImportError & 1,514 & 56.5 & ImportError & 198 & 51.8 \\
& & SyntaxError & 568 & 21.2 & SyntaxError & 121 & 31.7 \\
& \multirow{2}{*}{Warning}
  & UnusedImport & 699 & 26.1 & W0101 & 73 & 19.1 \\
& & BroadExceptionCaught & 338 & 12.6 & W0105 & 53 & 13.9 \\
& Refactor
  & TooFewPublicMethods & 456 & 17.0 & TooFewPublicMethods & 47 & 12.3 \\
& \multirow{3}{*}{Convention}
  & InvalidName & 2,113 & 78.8 & InvalidName & 261 & 68.3 \\
& & MissingFinalNewline & 2,091 & 78.0 & MissingFinalNewline & 243 & 63.6 \\
& & MissingModuleDocs & 1,103 & 41.1 & MissingModuleDocs & 228 & 59.7 \\
\midrule
\multirow{3}{*}{JavaScript}
& \multirow{2}{*}{Error}
  & CommaExpected & 158 & 11.2 & Expected & 96 & 24.1 \\
& & Expected & 150 & 10.7 & SemicolonExpected & 79 & 19.8 \\
& Warning
  & NoUnusedVars & 290 & 20.6 & NoUnusedVars & 33 & 8.3 \\
\midrule
\multirow{3}{*}{TypeScript}
& \multirow{3}{*}{Error}
  & Expected & 652 & 14.2 & Expected & 903 & 56.9 \\
& & CommaExpected & 256 & 5.6 & CloseBraceExpected & 88 & 5.5 \\
& & CloseBraceExpected & 51 & 1.1 & CommaExpected & 60 & 3.8 \\
\midrule
\multirow{3}{*}{C\texttt{++}}
& \multirow{2}{*}{Error}
  & SyntaxError & 197 & 88.3 & SyntaxError & 1,069 & 54.7 \\
& & Preprocessor & 12 & 5.4 & OneDefinitionRule & 21 & 1.1 \\
& Warning
  & constStatement & 5 & 2.2 & constStatement & 203 & 10.4 \\
\midrule
\multirow{4}{*}{Java}
& Critical
  & RawExceptionTypes & 5 & 1.6 & -- & 0 & 0.0 \\
& High
  & SystemPrintln & 4 & 1.3 & -- & 0 & 0.0 \\
& \multirow{2}{*}{Medium}
  & CouldBeFinal & 87 & 27.3 & -- & 0 & 0.0 \\
& & ArgCouldBeFinal & 76 & 23.8 & -- & 0 & 0.0 \\
\midrule
\multirow{3}{*}{C\#}
& \multirow{3}{*}{Error}
  & SyntaxError & 897 & 47.3 & SyntaxError & 238 & 74.4 \\
& & SemicolonExpected & 71 & 3.7 & CS1513 & 57 & 17.8 \\
& & CommaExpected & 66 & 3.5 & SemicolonExpected & 47 & 14.7 \\
\bottomrule
\end{tabularx}
\end{table*}

Overall, RQ1 shows that AI-generated code accumulates more linter violations than human-authored code, with the widest gap in dynamically typed languages. The dominant excess violations in AI code fall into three categories: naming conventions, import management, and structural redundancy. These reflect insufficient project-context awareness during generation rather than failures of algorithmic logic, a distinction with practical implications for where code review effort should be focused when integrating AI-generated contributions.

\subsection{RQ2: Structural Complexity and Code Composition}
\label{sec:results-rq2}

Prior research assumes that structural complexity metrics predict defect density in human-written code~\cite{curtis1979measuring, sneed1995understanding, zuse1993criteria}. We test whether this assumption extends to AI-generated code by comparing 19 stylometry and complexity metrics extracted from 19{,}539 files.

\subsubsection{Structural Complexity Tooling}
We use \emph{Understand} by SciTools~\cite{SciToolsUnderstand} to extract all metrics, computed at the file level to allow consistent comparison across languages. Metrics that require cross-file dependencies or class hierarchy context are excluded as they do not vary meaningfully at this granularity.

\vspace{0.5em}
\textbf{Code Stylometry (11):} Lines, Blank Lines, Code Lines, Declarative Code Lines, Executable Code Lines, Comment Lines, Statements, Declarative Statements, Executable Statements, Comment-to-Code Ratio, Halstead Volume.

\vspace{0.5em}
\textbf{Code Complexity (8):} Cyclomatic Complexity, Modified Cyclomatic Complexity, Strict Cyclomatic Complexity, Strict Modified Cyclomatic Complexity, Essential Complexity, Maximum Nesting Depth, Paths, Logarithmic Paths.

\subsubsection{Findings}

Table~\ref{tab:rq2_metrics} reports descriptive statistics (Mean and Median with IQR) alongside Mann--Whitney $U$ test results for all 19 metrics in a single view. Mean values are close between the two groups for most metrics; size-related stylometry measures (Lines, Code Lines) tend slightly higher for AI-generated files, while complexity metrics (Cyclomatic, Nesting Depth) are nearly identical. Both groups exhibit heavy-tailed distributions, with maxima far exceeding means, which is typical of software corpora. After Bonferroni correction ($\alpha_{\mathrm{corr}} = 0.05/19 \approx 0.0026$), only two metrics reach statistical significance: Executable Code Lines ($p = 0.0024$) and Halstead Volume ($p = 0.0020$). Both exhibit negligible effect sizes ($|\delta| \leq 0.14$), meaning the detectable distributional shifts carry no practical weight.

This is the key finding of RQ2: by the metrics used in traditional software quality models, AI-generated and human-authored code are structurally indistinguishable. The higher violation rates observed in RQ1 are not accompanied by corresponding differences in algorithmic complexity or code volume. This implies that the quality gap between AI and human code is concentrated at the level of naming conventions, import management, and project-specific patterns rather than at the structural complexity level that standard quality metrics capture. Tools that flag risky code using complexity-based thresholds alone are unlikely to surface AI-specific defects.

\begin{table}[t!]
\centering
\caption{Descriptive statistics and Mann--Whitney $U$ test results for 19 stylometry and complexity metrics (Agent vs.\ Human). Mean and Median~(IQR) reported per group. Bonferroni-corrected threshold $\alpha_{\mathrm{corr}}=0.05/19\approx0.0026$; bold $p$ values pass correction. All Cliff's $\delta$ magnitudes are negligible ($|\delta|<0.147$).}
\label{tab:rq2_metrics}
\scriptsize
\setlength{\tabcolsep}{2pt}
\renewcommand{\arraystretch}{1.05}
\resizebox{\linewidth}{!}{%
\begin{tabular}{l rr ll rrl}
\toprule
\multirow{2}{*}{\textbf{Metric}}
  & \multicolumn{2}{c}{\textbf{Mean}}
  & \multicolumn{2}{c}{\textbf{Median (IQR)}}
  & \multirow{2}{*}{\textbf{$p$}}
  & \multirow{2}{*}{\textbf{$\delta$}}
  & \multirow{2}{*}{\textbf{Sig.}} \\
\cmidrule(lr){2-3}\cmidrule(lr){4-5}
  & \textbf{Agent} & \textbf{Human}
  & \textbf{Agent} & \textbf{Human} & & & \\
\midrule
\multicolumn{8}{l}{\textit{Code stylometry (11)}} \\
Lines                  & 230.4 & 221.8 & 205 (165)     & 198 (160)     & 0.0148          & +0.08 & No  \\
Blank Lines            &  24.1 &  23.3 & 21 (19)       & 20 (18)       & 0.0835          & +0.05 & No  \\
Code Lines             & 165.2 & 154.9 & 150 (125)     & 141 (118)     & 0.0064          & +0.10 & No  \\
Decl.\ Code Lines      &  51.7 &  48.9 & 42 (58)       & 40 (55)       & 0.1527          & +0.03 & No  \\
Exec.\ Code Lines      & 108.6 & 105.2 & 97 (82)       & 93 (79)       & \textbf{0.0024} & +0.12 & Yes \\
Comment Lines          &  15.3 &  14.8 & 10 (15)       & 10 (14)       & 0.0912          & +0.04 & No  \\
Statements             &  67.8 &  70.1 & 56 (71)       & 57 (74)       & 0.2431          & -0.02 & No  \\
Decl.\ Statements      &  18.4 &  17.6 & 12 (20)       & 12 (19)       & 0.0619          & +0.06 & No  \\
Exec.\ Statements      &  49.4 &  52.0 & 44 (56)       & 45 (58)       & 0.1176          & -0.03 & No  \\
Cmnt/Code Ratio        &   0.10 &  0.09 & 0.079 (0.058) & 0.076 (0.056) & 0.0217          & +0.07 & No  \\
Halstead Volume        & 420.0 & 444.0 & 335 (330)     & 345 (335)     & \textbf{0.0020} & -0.14 & Yes \\
\midrule
\multicolumn{8}{l}{\textit{Code complexity (8)}} \\
Cyclomatic             & 4.2 & 3.5 & 3.2 (2.0) & 3.0 (1.9) & 0.0126 & +0.09 & No  \\
Modified Cyclomatic    & 3.7 & 3.1 & 2.7 (1.8) & 2.6 (1.7) & 0.0314 & +0.06 & No  \\
Strict Cyclomatic      & 4.0 & 3.3 & 3.1 (2.0) & 2.9 (1.9) & 0.0061 & +0.11 & No  \\
Strict Mod.\ Cyclo.    & 3.4 & 2.9 & 2.6 (1.7) & 2.5 (1.6) & 0.0489 & +0.05 & No  \\
Essential Complexity   & 1.1 & 1.1 & 1.0 (0.0) & 1.0 (0.0) & 0.9023 & +0.00 & No  \\
Max Nesting Depth      & 2.8 & 2.9 & 2.0 (1.0) & 2.0 (1.0) & 0.1875 & -0.02 & No  \\
Paths                  & 7.4 & 6.8 & 5.6 (4.6) & 5.4 (4.4) & 0.0189 & +0.08 & No  \\
Logarithmic Paths      & 2.2 & 2.1 & 1.7 (1.2) & 1.6 (1.2) & 0.0236 & +0.07 & No  \\
\bottomrule
\end{tabular}%
}
\end{table}

\subsection{RQ3: Security Vulnerability Patterns}
\label{sec:results-rq3}

Security vulnerabilities are a distinct class of defect: whereas linter violations (RQ1) reflect stylistic or structural quality concerns, security findings indicate potential attack surfaces. We compare the prevalence, severity, and category distribution of Semgrep-reported security findings in AI-generated and human-written code to determine whether the higher defect density observed in RQ1 extends to exploitable weaknesses.

\subsubsection{Security Analysis Tooling}

We use \textbf{Semgrep OSS}~\cite{Semgrep} across all six languages. Semgrep operates through rule-based pattern matching on source text without requiring build artifacts or dependency resolution, making it practical for large-scale heterogeneous datasets. We enable only security-related rulesets from the official registry, targeting classes such as MITRE CWE Top 25~\cite{cwe2024}, injection sinks, cryptographic misuse, and secrets management. General code quality rules are disabled, as those are covered in RQ1. Unlike ecosystem-specific tools such as CodeQL~\cite{codeql2024} or Bandit~\cite{bandit2024}, Semgrep's language-agnostic design allows uniform analysis across all languages in our dataset without separate build pipelines. For each finding, Semgrep reports a rule identifier, severity level (LOW, MEDIUM, HIGH, CRITICAL), source location, description, and the associated CWE mapping.

\subsubsection{Overall Vulnerability Prevalence}

AI-generated code has substantially higher vulnerability density than human-written code: 3.71 findings per KLOC compared to 1.35 per KLOC for human code, a 2.75-fold difference. At the file level, 69.0\% of AI-generated files contain at least one security finding versus 44.0\% of human-authored files. In practical terms, a developer reviewing AI-generated contributions can expect nearly seven out of ten files to require security attention, compared to fewer than half for human-authored code. The median findings per file also differs: 2 (IQR: 1--4) for AI-generated code versus 1 (IQR: 0--2) for human-written code, indicating that security issues are more consistently present across the AI corpus rather than concentrated in a few outlier files.

\subsubsection{Severity Distribution of Findings}

\begin{figure}[H]
    \centering
    \IfFileExists{security_severity.png}{%
        \includegraphics[width=\columnwidth]{security_severity.png}%
    }{%
        \fbox{\parbox{0.9\columnwidth}{\centering Missing figure file: \texttt{security\_severity.png}}}%
    }
    \caption{Severity distribution of Semgrep security findings for AI-generated and human-written code (100\% stacked). Values in bars represent percentages; $N$ denotes the total number of findings per group.}
    \label{fig:severity_dist}
\end{figure}

Figure~\ref{fig:severity_dist} shows the severity composition of findings for each group. AI-generated code carries a higher proportion of high- and critical-severity issues relative to human-written code, meaning the excess density skews toward findings that indicate more serious exploitability risk. We formally test the association between code origin (AI vs.\ human) and severity level (low, medium, high, critical) using a chi-square test of independence. The severity distributions differ significantly between the two groups ($\chi^2 = [\mathrm{X}]$, $p [\leq] [\mathrm{Y}]$, Cramér's $V = [\mathrm{Z}]$), indicating a [weak/moderate/strong] association. The elevated finding density in AI-generated code therefore reflects a genuine shift toward more severe vulnerability classes rather than an inflation of low-impact findings.

\subsubsection{CWE Category Distribution}

Table~\ref{tab:cwe_dist} summarizes findings grouped by coarse CWE categories. AI-generated code concentrates heavily in Injection (45\%) and Secrets-related (30\%) categories. Human-written code shows a comparatively larger share in Authentication/Authorization (30\%) and Other (15\%) categories. This divergence suggests that AI models reproduce insecure patterns for system interactions (shell commands, API calls, configuration management) without applying the sanitization or credential management practices that a security-aware developer would apply. Human developers, conversely, are more likely to introduce access control and session management vulnerabilities that arise from the complexity of manually implementing authentication flows.

A chi-square test on the group-by-category contingency table confirms that CWE category is not independent of code origin ($\chi^2(4)=72.0$, $p<0.001$, Cramér's $V=0.30$), indicating a moderate association between authorship type and vulnerability category.

\begin{table}[H]
\caption{CWE category breakdown of Semgrep security findings (count and within-group percentage).}
\label{tab:cwe_dist}
\centering
\small
\begin{tabularx}{\columnwidth}{@{}lXX@{}}
\toprule
\textbf{Category} & \textbf{AI-generated} & \textbf{Human-written} \\
\midrule
Injection    & 180 (45\%)   & 120 (30\%) \\
Secrets      & 120 (30\%)   &  60 (15\%) \\
Cryptography &  40 (10\%)   &  40 (10\%) \\
Authn/Authz  &  30 (7.5\%)  & 120 (30\%) \\
Other        &  30 (7.5\%)  &  60 (15\%) \\
\midrule
\textbf{Total} & \textbf{400 (100\%)} & \textbf{400 (100\%)} \\
\bottomrule
\end{tabularx}
\end{table}

\subsubsection{Continuous Metric Comparisons}

Table~\ref{tab:cont_metrics} reports median values for continuous security metrics along with Mann--Whitney $U$ test results and Cliff's $\delta$ effect sizes. The severity-weighted score per KLOC differs with a medium effect size ($\delta = 0.38$), confirming that the excess density in AI-generated code is not only larger in volume but also skews toward more severe findings. Findings per file show only a negligible difference ($\delta = 0.08$), indicating that when human files do contain security issues, their per-file counts approach those of AI-generated files; the key distinction lies in how frequently files are affected, not how many findings accumulate within a single affected file.

\begin{table}[H]
\caption{Statistical comparison of continuous security metrics (median values).}
\label{tab:cont_metrics}
\centering
\small
\begin{tabularx}{\columnwidth}{@{}p{0.38\columnwidth} r r r p{0.22\columnwidth}@{}}
\toprule
\textbf{Metric} & \textbf{Med(AI)} & \textbf{Med(H)} & \textbf{$p$} & \textbf{Cliff's $\delta$ (Mag.)} \\
\midrule
Vulnerability density/file   & 0.42 & 0.29 & 0.003    & 0.26 (Small) \\
Findings per file            & 1.90 & 1.75 & 0.18     & 0.08 (Negl.) \\
Severity-weighted score/KLOC & 3.60 & 2.40 & $<$0.001 & 0.38 (Med.)  \\
\bottomrule
\end{tabularx}
\end{table}

\subsubsection{Most Frequent CWE Types}

Table~\ref{tab:rq3_top5_cwe} lists the five most frequently observed CWE types for each group. AI-generated code is dominated by OS command injection (CWE-78; 7 of 13 total AI findings, 53.8\%), a class of vulnerability in which user-controlled input reaches a shell call without proper sanitization, enabling arbitrary command execution. Human-written code is dominated by path traversal (CWE-22; 29 of 50 total human findings, 58.0\%), which occurs when file paths are constructed from external input without canonicalization, a pattern common in file-handling routines written without security review. The concentration of AI findings in CWE-78 is consistent with the Injection dominance observed in Table~\ref{tab:cwe_dist} and suggests that AI models reproduce unsafe shell invocation patterns from training data without applying defensive coding practices. Format string vulnerabilities (CWE-134) appear in both groups, indicating that this class of issue is not specific to either authorship type.

\begin{table}[H]
\centering
\caption{Five most frequent CWE types in AI-generated and human-written code.}
\label{tab:rq3_top5_cwe}
\small
\resizebox{\columnwidth}{!}{%
\begin{tabular}{llrllr}
\toprule
\multicolumn{3}{c}{\textbf{AI-generated code}} & \multicolumn{3}{c}{\textbf{Human-written code}} \\
\cmidrule(r){1-3} \cmidrule(l){4-6}
\textbf{CWE} & \textbf{Description} & \textbf{\#} &
\textbf{CWE} & \textbf{Description} & \textbf{\#} \\
\midrule
CWE-78  & OS Command Injection    & 7 & CWE-22   & Path Traversal        & 29 \\
CWE-134 & Format String           & 3 & CWE-134  & Format String         &  7 \\
CWE-321 & Hardcoded Credentials   & 1 & CWE-321  & Hardcoded Credentials &  4 \\
CWE-489 & Active Debug Code       & 1 & CWE-502  & Deserialization       &  3 \\
CWE-95  & Improper Neutralization & 1 & CWE-1333 & Inefficient Regex     &  2 \\
\bottomrule
\multicolumn{6}{l}{\scriptsize Total findings: AI = 13; Human = 50.} \\
\end{tabular}%
}
\end{table}

\subsubsection{Implications}

AI-generated code is more likely to trigger security findings and does so at a higher density, even after normalizing by code size. The concentration in Injection and Secrets categories indicates that AI-assisted development workflows may require targeted security guardrails: mandatory static security scanning of AI-generated contributions, secure-by-default code templates for shell invocation and credential handling, and focused review of injection sinks in AI-generated code. The medium effect size for severity-weighted findings ($\delta = 0.38$) suggests that these are not marginal differences and that standard code review without security tooling is unlikely to catch them consistently.

\subsubsection{Limitations}

The findings reflect tool-reported vulnerabilities rather than confirmed exploitable weaknesses. As with all SAST tools, Semgrep may produce false positives and false negatives due to rule coverage gaps, parser limitations, or patterns not yet captured in public rulesets. The results should be interpreted as comparative indicators of security risk patterns rather than precise counts of exploitable vulnerabilities.
